% ============================================================
%  chapter-template.tex — Markdown-to-LaTeX Conversion Guide
%  Math for Machine Learning: A Software Engineer's Guide
%
%  This file is NOT included in book.tex.  It serves as a
%  reference showing how each Markdown construct maps to its
%  LaTeX equivalent when converting chapters by hand or when
%  reviewing pandoc output.
%
%  Pandoc one-liner (run from the repo root):
%    pandoc chapters/07-probability-foundations.md \
%      --from markdown --to latex \
%      --listings \
%      -o latex/ch07.tex
% ============================================================

% ------------------------------------------------------------
% HEADING MAPPING
% ------------------------------------------------------------
%
%  Markdown                          LaTeX
%  ────────────────────────────────  ─────────────────────────
%  # Chapter N: Title                \chapter{Title}
%  ## 9.2 Section Title              \section{Section Title}
%  ### 9.2.1 Subsection Title        \subsection{Subsection Title}
%  #### Sub-subsection Title         \subsubsection{Sub-subsection Title}
%
% Example:
\chapter{Probability Foundations}         % from # Chapter 7: Probability Foundations
\section{Sample Spaces and Events}        % from ## 7.1 Sample Spaces and Events
\subsection{The Kolmogorov Axioms}        % from ### 7.1.1 The Kolmogorov Axioms

% ------------------------------------------------------------
% INLINE MATH
% ------------------------------------------------------------
%
%  Markdown: The probability $P(A) \in [0, 1]$ for any event $A$.
%  LaTeX:    The probability $P(A) \in [0, 1]$ for any event $A$.
%            (identical — pandoc preserves $ delimiters)

The probability $P(A) \in [0, 1]$ for any event $A$.

% ------------------------------------------------------------
% DISPLAY MATH (equation environment)
% ------------------------------------------------------------
%
%  Markdown:
%    $$
%    \mathcal{L}(\theta) = \prod_{i=1}^{n} p(x_i \mid \theta)
%    $$
%
%  LaTeX (pandoc output uses equation*; use equation for numbered):

\begin{equation}
  \mathcal{L}(\theta) = \prod_{i=1}^{n} p(x_i \mid \theta)
  \label{eq:likelihood}
\end{equation}

% Cross-reference: see equation~\eqref{eq:likelihood}.

% Multi-line aligned equations:
\begin{align}
  \ell(\theta)
    &= \log \mathcal{L}(\theta) \notag \\
    &= \sum_{i=1}^{n} \log p(x_i \mid \theta)
    \label{eq:loglik}
\end{align}

% ------------------------------------------------------------
% BOLD / ITALIC TEXT
% ------------------------------------------------------------
%
%  Markdown: **bold**, *italic*, ***bold italic***
%  LaTeX:    \textbf{bold}, \textit{italic}, \textbf{\textit{bold italic}}

\textbf{Maximum Likelihood Estimation} (MLE) is the cornerstone of
\textit{parametric} statistics.

% ------------------------------------------------------------
% BLOCK QUOTE
% ------------------------------------------------------------
%
%  Markdown:
%    > "A model that trains well is promising. A model you can trust is useful."
%
%  LaTeX:

\begin{quote}
  \itshape
  ``A model that trains well is promising.
    A model you can trust is useful.''
\end{quote}

% ------------------------------------------------------------
% UNORDERED LIST
% ------------------------------------------------------------
%
%  Markdown:
%    - First item
%    - Second item
%      - Nested item
%
%  LaTeX:

\begin{itemize}
  \item First item
  \item Second item
    \begin{itemize}
      \item Nested item
    \end{itemize}
\end{itemize}

% ------------------------------------------------------------
% ORDERED LIST
% ------------------------------------------------------------
%
%  Markdown:
%    1. Compute the likelihood.
%    2. Take the log.
%    3. Differentiate and set to zero.
%
%  LaTeX:

\begin{enumerate}
  \item Compute the likelihood.
  \item Take the log.
  \item Differentiate and set to zero.
\end{enumerate}

% ------------------------------------------------------------
% TABLE
% ------------------------------------------------------------
%
%  Markdown:
%    | Fold | Val indices | Fold MSE |
%    |------|-------------|----------|
%    | 1    | 0, 1        | 20.141   |
%    | 2    | 2, 3        | 8.500    |
%
%  LaTeX (booktabs style):

\begin{table}[htbp]
  \centering
  \caption{5-Fold cross-validation results.}
  \label{tab:cv-results}
  \begin{tabular}{ccc}
    \toprule
    Fold & Val indices & Fold MSE \\
    \midrule
    1 & 0, 1 & 20.141 \\
    2 & 2, 3 & 8.500  \\
    \bottomrule
  \end{tabular}
\end{table}

% ------------------------------------------------------------
% CODE LISTING (fenced code block)
% ------------------------------------------------------------
%
%  Markdown:
%    ```python
%    import numpy as np
%    w = np.linalg.solve(X.T @ X, X.T @ y)
%    print(w)
%    ```
%
%  LaTeX (listings package, style defined in preamble.tex):

\begin{lstlisting}[style=python, caption={Normal Equation solution.}, label={lst:normal-eq}]
import numpy as np

# Normal Equation: w = (X^T X)^{-1} X^T y
w = np.linalg.solve(X.T @ X, X.T @ y)
print(w)
\end{lstlisting}

% Plain/output listing (no syntax highlighting):
\begin{lstlisting}[style=text, caption={Program output.}]
Weights: [1.234  0.567  -0.890]
\end{lstlisting}

% ------------------------------------------------------------
% FIGURE
% ------------------------------------------------------------
%
%  Markdown: (no native figure syntax; images use ![caption](path))
%
%  LaTeX:

\begin{figure}[htbp]
  \centering
  \includegraphics[width=0.75\textwidth]{figures/bias-variance-tradeoff}
  \caption{The bias-variance tradeoff: model complexity vs.\ test error.}
  \label{fig:bias-variance}
\end{figure}

% Cross-reference: Figure~\ref{fig:bias-variance} shows \ldots

% ------------------------------------------------------------
% DEFINITION / THEOREM ENVIRONMENTS
% ------------------------------------------------------------
%
%  These environments are declared in preamble.tex.

\begin{definition}[Sample Space]
  The \textbf{sample space} $\Omega$ is the set of all possible outcomes
  of a random experiment.
\end{definition}

\begin{theorem}[Bayes' Theorem]
  For events $A$ and $B$ with $P(B) > 0$:
  \begin{equation}
    P(A \mid B) = \frac{P(B \mid A)\,P(A)}{P(B)}.
  \end{equation}
\end{theorem}

\begin{example}[Medical Test]
  A test has sensitivity 99\% and specificity 99\%.
  If the disease prevalence is 1 in 10{,}000, then by Bayes' theorem
  the positive predictive value is approximately 1\%.
\end{example}

% ------------------------------------------------------------
% CROSS-REFERENCES
% ------------------------------------------------------------
%
%  Always use \label{} on equations, figures, tables, and sections.
%  Reference with \ref{} (number only) or \eqref{} (equation with parens).

% \section{The Bias-Variance Tradeoff}
% \label{sec:bias-variance}
%
% As derived in equation~\eqref{eq:loglik} and illustrated in
% Figure~\ref{fig:bias-variance}, the decomposition holds for any
% squared-error loss (see also Table~\ref{tab:cv-results}).
